%! AP Statistics — Unit 7, Lesson 7.3 — Warm-Up (Answer Key)
\documentclass[10pt]{article}
\usepackage{apstats-article}
\usepackage{apstats-key}
\newcommand{\TallMath}[1]{$\displaystyle #1\rule[-1.4em]{0pt}{3.2em}$}

\begin{document}

\pageheader{Unit 7, Lesson 7.3}{Warm-Up --- Answer Key}
\namedateperiod

\vspace{0.15in}

\textbf{95\% CI for mean sodium intake: $(2615, 3065)$ mg}

\begin{enumerate}

  \item Three interpretations:

        \begin{enumerate}[label=\Alph*.]
          \item ``There is a 95\% probability that the true mean sodium intake is
                between 2615 and 3065 mg.''

                Error: \ansline{Uses ``probability'' --- implies $\mu$ is random. $\mu$ is fixed; it is either in the interval or not. Use ``confident'' instead.}

          \item ``We are 95\% confident that the interval (2615, 3065) mg captures
                the true mean daily sodium intake for all college students.''

                Error: \ans{None --- this is the correct interpretation.} \textcolor{keyred}{\textbf{$\leftarrow$ correct}}

          \item ``95\% of all college students have a daily sodium intake between
                2615 and 3065 mg.''

                Error: \ansline{Confuses a CI for the parameter $\mu$ with a range for individual data values. A CI estimates the mean, not individual observations.}
        \end{enumerate}

  \item Does $(0.546, 0.734)$ provide convincing evidence against $p = 0.80$?
        \ansline{Yes: the claimed value 0.80 lies outside the interval (0.546, 0.734), so the data provide convincing evidence that the true proportion differs from 0.80.}

  \item Effect of doubling $n$:
        \ansline{The margin of error would decrease by a factor of $1/\sqrt{2} \approx 0.707$ --- approximately 29\% smaller --- because $ME \propto 1/\sqrt{n}$ and $\sqrt{40}/\sqrt{20} = \sqrt{2}$.}

\end{enumerate}

\end{document}
